\documentclass[aspectratio=169,10pt]{beamer}
\usetheme[numbering=fraction,progressbar=frametitle]{metropolis}
\usepackage{fontspec}
\setsansfont{Latin Modern Sans}
\usepackage{amsmath,amssymb,mathtools,graphicx,booktabs}
\usepackage{polyglossia}
\setdefaultlanguage{english}
\setbeamertemplate{navigation symbols}{}
\setbeamerfont{frametitle}{size=\large}
\setbeamerfont{normal text}{size=\normalsize}
\setlength{\parskip}{.45em}
\hypersetup{pdfauthor={},pdfsubject={Lecture slides},pdfkeywords={},pdftitle={The logarithm function}}
\providecommand{\N}{\mathbb{N}}
\providecommand{\Z}{\mathbb{Z}}
\providecommand{\D}{\mathbb{D}}
\providecommand{\Q}{\mathbb{Q}}
\providecommand{\R}{\mathbb{R}}
\providecommand{\C}{\mathbb{C}}
\providecommand{\vect}[1]{\overrightarrow{#1}}
\providecommand{\vi}{\vec{\imath}}
\providecommand{\vj}{\vec{\jmath}}
\providecommand{\vk}{\vec{k}}
\providecommand{\norm}[1]{\left\lVert #1\right\rVert}
\providecommand{\vc}[2]{\begin{pmatrix}#1\\#2\end{pmatrix}}
\providecommand{\deter}[2]{\det\left(#1,#2\right)}
\providecommand{\abs}[1]{\left\lvert #1\right\rvert}
\providecommand{\dd}{\,\mathrm{d}}
\providecommand{\e}{\mathrm{e}}
\providecommand{\Card}{\operatorname{Card}}
\providecommand{\E}{\operatorname{E}}
\providecommand{\V}{\operatorname{V}}
\providecommand{\degre}{^{\circ}}
\providecommand{\un}[1]{\,\mathrm{#1}}
\providecommand{\sci}[2]{#1\times10^{#2}}
\providecommand{\pgcd}{\operatorname{PGCD}}
\providecommand{\ppcm}{\operatorname{PPCM}}
\title{The logarithm function}
\subtitle{Lecture and worked examples}
\date{}
\author{}
\begin{document}
\begin{frame}[plain]
\vfill
{\usebeamerfont{title}\inserttitle\par}
\medskip
{\usebeamerfont{subtitle}\insertsubtitle\par}
\medskip
{\usebeamercolor[fg]{progress bar}\rule{\textwidth}{.5pt}}
\vfill
\end{frame}
\begin{frame}{Objectives}
Use the properties of ln; solve equations and thresholds; differentiate and study logarithmic functions.
\end{frame}
\begin{frame}{Definition and properties}
The exponential is continuous and strictly increasing on \(\R\), and its range is \(]0,+\infty[\). For every \(x>0\), there is a unique real number \(y\) such that \(e^y=x\): it is denoted \(\ln x\).
\[\ln(e^y)=y\quad(y\in\R),\qquad e^{\ln x}=x\quad(x>0),\qquad\ln1=0,\quad\ln e=1.\]
For \(a,b>0\) and \(n\in\Z\),
\[\ln (ab)=\ln a+\ln b,\quad\ln (\dfrac{a}{b})=\ln a-\ln b,\quad\ln (a^n)=n\ln a,\quad\ln \sqrt{a}=\tfrac{1}{2}\ln a.\]
\end{frame}
\begin{frame}{Definition and properties}
\(e^{\ln a+\ln b}=e^{\ln a}e^{\ln b}=ab\); the uniqueness of the preimage gives the formula. There is no analogous formula for \(\ln(a+b)\).
\end{frame}
\begin{frame}{Derivative, variations and limits}
We accept that ln is differentiable on \(]0,+\infty[\). Differentiating \(e^{\ln x}=x\) gives \(e^{\ln x}(\ln x)'=1\); hence \((\ln x)'=\dfrac{1}{x}>0\). Ln is strictly increasing and concave, since its second derivative is \(-\dfrac{1}{x^2}\).
\[\lim_{x\to0^+}\ln x=-\infty,\qquad\lim_{x\to+\infty}\ln x=+\infty.\]
Its curve is the reflection of that of the exponential in the line \(y=x\). The line \(x=0\) is a vertical asymptote. Its tangent at \(1\) gives \(\ln x\le x-1\) for \(x>0\).
\end{frame}
\begin{frame}{Derivative, variations and limits}
If \(u>0\) is differentiable, \((\ln u)'=\dfrac{u'}{u}\). For every integer \(m\ge1\),
\[\lim_{x\to+\infty}\dfrac{\ln x}{x^m}=0,\qquad\lim_{x\to0^+}x^m\ln x=0.\]
\end{frame}
\begin{frame}{Derivative, variations and limits}
Set \(t=-\ln x\); as \(x\to0^+\), \(t\to+\infty\) and \(x\ln x=-te^{-t}\to0\) by comparative growth. For \(m\ge1\), the same reasoning gives \(-te^{-mt}\to0\). At \(+\infty\), set \(t=\ln x\) to write \(\dfrac{\ln x}{x^m}=\dfrac{t}{e^{mt}}\to0\).
\end{frame}
\begin{frame}{Equations and thresholds}
For \(A,B>0\), \(\ln A=\ln B\Leftrightarrow A=B\) and \(\ln A\le\ln B\Leftrightarrow A\le B\). Always start with the domain.
\end{frame}
\begin{frame}{Equations and thresholds}
\(\ln(x-1)+\ln(x+1)=\ln8\) requires \(x>1\). The equation becomes \(x^2-1=8\), so \(x=3\) only.
\end{frame}
\begin{frame}{Equations and thresholds}
\(0.8^n\le0.1\) is equivalent to \(n\ln0.8\le\ln0.1\), then \(n\ge\dfrac{\ln 0.1}{\ln 0.8}\simeq10.319\). The first integer is \(11\).
\end{frame}
\begin{frame}{Equations and thresholds}
\begin{columns}[T,onlytextwidth]
\begin{column}{.57\textwidth}
For \(f(x)=\dfrac{\ln x}{x}\), \(f'=\dfrac{1-\ln x}{x^2}\); \(f\) increases on \(]0,e]\), decreases on \([e,+\infty[\), and its maximum is \(\dfrac{1}{e}\). Its limit at \(0^+\) is \(-\infty\), its limit at \(+\infty\) is \(0\).
\end{column}
\begin{column}{.4\textwidth}
\centering
\includegraphics[width=\linewidth,height=4.5cm,keepaspectratio]{figures/terminale-08-1-en.pdf}
\end{column}
\end{columns}
\end{frame}
\begin{frame}{Research activity}
\textbf{Controlling a calculation error - 30 min.}

A calculation method produces an initial error \(e_0=5\). At each step, the error is multiplied by \(0.7\). Let \(e_n\) denote the error after \(n\) steps, for \(n\in\N\).
\end{frame}
\begin{frame}{Research activity (2)}
\begin{enumerate}
\setcounter{enumi}{0}\item Express \(e_n\) in terms of \(n\).

\item Using the logarithm, find the smallest integer \(n\) such that \(e_n\le10^{-4}\).

\item Using a calculator, compute the error at the value found and at the previous one to check that it is the minimum.

\item Now require an error of at most \(10^{-6}\). Find the new minimum number of steps. Does dividing the tolerance by \(100\) multiply the number of steps by \(100\)?

\end{enumerate}
\end{frame}
\end{document}
